\input{../../../stock/en-tete_v6.tex}

\newcommand{\titrechapitre}{Fonctions}
\newcommand{\numerochapitre}{1}

\renewcommand{\headrulewidth}{0.4pt}
\renewcommand{\footrulewidth}{0.4pt}
\fancyfoot[L]{Raphaël JONTEF}
\fancyfoot[C]{\thepage}
\fancyfoot[R]{\href{https://www.alveusclub.com/}{Alveus}}
\fancyhead[L]{Mathématiques : Troisième}
\fancyhead[R]{\titrechapitre}

\begin{document}
\input{../../../stock/commands.tex}

% Page de titre custom
\setlength{\headheight}{15pt}
\begin{titlepage}
    \centering
    \vspace*{3cm}
    {\huge Chapitre \numerochapitre\par}
    {\Huge \bfseries\titrechapitre\par}
    \vspace{8cm}
    {\Large Cours rédigé par Raphaël JONTEF\par}
    \vspace{2em}
    {\small \textit{raphaeljontef@hotmail.com}\par}
    \vspace{2cm}

    {\normalsize pour un enseignement dans le cadre  \par}
    {\small d'un soutien scolaire encadré par \href{https://www.alveusclub.com/}{Alveus} \par}
    \vfill
    {\large \today\par}
\end{titlepage}

% plan du cours
\tableofcontents
\newpage


% -----------Corps du document--------------

\section{Notion de fonction}
\subsection{Définitions et vocabulaire}
Une fonction est une relation qui associe à chaque élément d'un ensemble de départ (appelé \notion{ensemble de départ}) un unique élément d'un ensemble d'arrivée (appelé \notion{ensemble d'arrivée}). On note $f : A \to B$ une fonction \(f\) qui associe chaque élément de l'ensemble \(A\) à un élément de l'ensemble \(B\).

\begin{definition}{}{antécédent et image}
    Soit \(f : A \to B\) une fonction.
    \begin{itemize}
        \item Pour un élément \(x\) de l'ensemble de départ \(A\), l'élément \(f(x)\) de l'ensemble d'arrivée \(B\) est appelé \notion{image} de \(x\) par la fonction \(f\).
        \item Pour un élément \(y\) de l'ensemble d'arrivée \(B\), un élément \(x\) de l'ensemble de départ \(A\) tel que \(f(x) = y\) est appelé \notion{antécédent} de \(y\) par la fonction \(f\).
    \end{itemize}
\end{definition}

\begin{exemple}{}{de fonction}
    \fonction{f}{\N}{\N}{n}{2n}
    est une fonction qui à chaque entier naturel \(n\) associe l'entier naturel \(2n\). Par exemple, l'image de \(3\) par la fonction \(f\) est \(f(3) = 6\). \\
    Moralement, la fonction \(f\) "double" chaque entier naturel. \\
\end{exemple}

\begin{exemple}{}{de fonction plus étrange}
    Il existe aussi des fonctions définies sur des ensembles quelconques. Par exemple, on peut définir la fonction suivante~:
    \fonction{f}{\{\text{Bernard}, \text{Claude}, \text{Michel}, \text{Patrick}\}}{\{\text{Claudette}, \text{Michèle}, \text{Patrick}, \text{}\}}{x}{\text{[le prénom féminin de $x$]}} \\
    On a alors~:
    \begin{itemize}
        \item $f(\text{Bernard}) = \text{Bernadette}$
        \item $f(\text{Claude}) = \text{Claudette}$
    \end{itemize}
    Mathématiquement, cette fonction est tout à fait déinie, même si elle n'a pas d'utilité mathématique.
\end{exemple}

\subsection{Fonctions affines}

On appelle \notion{fonction affine} toute fonction \(f : \R \to \R\) qui peut s'écrire sous la forme~:
$$f : x \mapsto ax + b$$
où \(a\) et \(b\) sont des nombres réels fixés.
\begin{itemize}
    \item Le nombre \(a\) est appelé le \notion{coefficient directeur} de la fonction affine.
    \item Le nombre \(b\) est appelé l'\notion{ordonnée à l'origine}.
\end{itemize}

\begin{exemple}{}{fonction affine}
    La fonction \(f : \R \to \R\) définie par \(f(x) = 2x + 3\) est une fonction affine de coefficient directeur \(2\) et d'ordonnée à l'origine \(3\).
\end{exemple}

\begin{exemple}{}{}
    L'antécédent de \(7\) par la fonction $f : \mapsto 2x + 3$ \\ est le nombre \(2\). En effet~:
    \begin{align*}
        f(2) & = 2 \times 2 + 3 \\
             & = 4 + 3 \\
             & = 7
    \end{align*}
\end{exemple}

\end{document}